\subsection{Finite support and semiperfect upper rings}\label{sec:finite-inverse}
We prove Theorem~\ref{thm:finite-inverse-projectives}. Its ordinary
semiperfect-ring input is the finite local-idempotent characterization
and projective-cover theory, not a finiteness assertion about the
coefficient groups; see Krause \cite[Proposition~4.1 and
Theorem~4.2]{Krause}.

\begin{proof}
\emph{Support and coefficient formulas.}
For fixed $l$, the admissible functionals form a nonempty finite set.
Their pointwise join preserves zero, finite joins and left $Q$-linearity
by distributivity. Its value at $k$, acting on $l$, is at most $k$,
so $\delta_l$ is admissible and greatest. No commutativity is used.
If $L$ is projective, choose a finite free retraction with coordinates
$\phi_i$ and generator images $l_i$. The identity
$k=\bigvee_i\phi_i(k)l_i$ makes $\phi_i$ admissible at $l_i$.
Replacing it by $\delta_{l_i}$ proves the support test, since each term
remains at most $k$. Conversely
\[
 \sigma(k)=(\delta_l(k))_{l\in L},\qquad
 \rho((q_l))=\bigvee_{l\in L}q_ll
\]
are $Q$-linear, and the test says exactly that $\rho\sigma=1_L$.

The set of $q$ with $ql\leq k$ is closed under joins and, since $Q$ is
finite, has greatest element $r(l,k)$. The coefficient-colimit formula
of Theorem~\ref{thm:action-colimits} therefore identifies its colimit
with $F(r(l,k))$. For an admissible $\phi$ at $l$,
$\phi(k)\leq\delta_l(k)$, and the coefficient transition gives
\[
 t_{\phi(k)l,k}a_{\phi(k),x,l}
 =t_{\delta_l(k)l,k}
   a_{\delta_l(k),F_{\phi(k),\delta_l(k)}(x),l}.
\]
Thus every canonical generator belongs to the displayed image from
$F(\delta_l(k))$. The reverse inclusion follows because $\delta_l$
is admissible. These formulas include $F(0)\ne0$.

\emph{Semiperfectness of the action ring.}
Write $e_l=r(l,l)$. Since $1l=l$, one has $1\leq e_l$. Also
$e_ll\leq l$, hence $e_l^2l\leq l$ and $e_l^2\leq e_l$ by
maximality. The reverse inequality follows from $1\leq e_l$, so
$e_l^2=e_l$. With $E_l$ denoting the identity block at the object $l$,
\[
 E_lB_LE_l=F(e_l)
\]
is exactly its ordinary unital coefficient ring. Its multiplication
needs no further support transition because $e_l^2=e_l$; its unit is
the image of the global unit of $T$.

The finitely many $E_l$ are orthogonal and sum to $1_{B_L}$.
Decompose each nonzero corner identity into a finite orthogonal family
of idempotents with local corner rings. If $f\leq E_l$ belongs to
such a family, then
\[
 fB_Lf=f(E_lB_LE_l)f,
\]
so its corner remains local in $B_L$. The combined families give a
finite orthogonal decomposition of $1_{B_L}$ into local idempotents.
The ordinary characterization makes $B_L$ semiperfect. A zero corner
contributes the empty family, and the zero ring is allowed.
No condition on off-diagonal finiteness or on radical preservation by
coefficient transitions has entered the argument.

Covariant additive $\mathcal A_L$-modules correspond to unital left
$B_L$-modules by taking the direct sum of the object groups; the inverse
takes the summands $E_lN$. This preserves finite generation and
projectivity, and identifies the actions of $\mathcal J_L$ and $I_L$.

\emph{Ordinary amplitudes and the ideal filter.}
Let $B$ be a semiperfect ring and let $J$ be its Jacobson radical.
Its simple left-module types $S_i$ have finitely generated
indecomposable projective covers $P_i$, with
$P_i/JP_i\simeq S_i$. Every finitely generated projective $P$ is a
finite sum of these covers, uniquely up to multiplicity. Here is the
usual argument for the last assertion. The top $P/JP$ is a finite
semisimple module. Form the finite sum $P'$ of the covers with those
simple multiplicities and lift its map to $P/JP$ through $P$.
The resulting $P'\to P$ is surjective by Nakayama's lemma and splits.
Its kernel $K$ is finitely generated, and equality of the tops gives
$K/JK=0$. Nakayama gives $K=0$. The multiplicities are then recovered
from the semisimple top. This uses no artinian hypothesis on $B$.

For any two-sided ideal $I$, $IS_i$ is either zero or all of $S_i$.
If $IP_i=P_i$, its image in $S_i$ is all. Conversely $IS_i\ne0$
gives
\[
 IP_i+\ker(P_i\twoheadrightarrow S_i)=P_i.
\]
The cover kernel is superfluous, so $IP_i=P_i$. Consequently a finite
sum $\bigoplus_iP_i^{n_i}$ satisfies $IP=P$ exactly when every present
type has $IS_i\ne0$; the ideal $I$ need not be idempotent or finitely
generated. Apply this with $(B,I)=(B_L,I_L)$ and use
Theorem~\ref{thm:canonical-ideal} to obtain precisely the stated
amplitudes. Zero amplitude retains the elements of $L$.

A support automorphism transports the action category and canonical
ideal, permuting the simple types and their covers. Semimodule
isomorphisms are exactly such support transports together with
amplitude isomorphisms. This gives the asserted orbit set. Finally,
a finitely generated projective $T$-semimodule is a retract of $T^n$,
so its support is a retract of the finite set $Q^n$. Thus every such
projective occurs in the classification.

\emph{The finite-semiring specialization.}
A finite unital ring is semiperfect. For completeness, repeatedly split
its finite regular module into indecomposable projectives. For any
endomorphism $f$ of a finite module, stabilization of kernels and images
gives $P=\ker(f^N)\oplus\im(f^N)$ for large $N$. On an
indecomposable summand, $f$ is therefore invertible or nilpotent;
if it is not invertible, $1-f$ is invertible. Its endomorphism ring
is consequently local. The summands of the regular module give the
required finite orthogonal local-idempotent decomposition.
Thus every finite inverse semiring satisfies the theorem's hypothesis.
Its arrow groups, and hence $B_L$, are finite. The formulas above use
only its finite operation and transition tables. They recover the
finite-semiring classification without imposing a hypothesis on
unitality of the individual lower coefficient fibers.
\end{proof}

\begin{example}[an infinite semiring with three support elements]\label{ex:finite-support-semiperfect}
Let $Q=\{\bot<1<u\}$, with addition maximum, absorbing $\bot$, and
multiplication maximum on the two nonbottom elements. For a field $K$,
take $F(\bot)=0$, $F(1)=F(u)=K$, with identity transition $F(1)\to F(u)$
and ordinary field products. These data define an inverse semiring:
all nonbottom coefficient transitions are identities and all coefficient
products are the field multiplication. Its global unit has support $1$.
Both upper coefficient rings are fields. If $K$ is infinite, the
semiring is infinite, while its support is finite and its unit is not
greatest.

For the regular support $L=Q$, order both the source columns and target
rows as $\bot,1,u$. Directly from $ql\leq k$,
\[
 (r(l,k))_{k,l}=
 \begin{pmatrix}u&\bot&\bot\\u&1&\bot\\u&u&u\end{pmatrix},
 \qquad
 B_L=\begin{pmatrix}K&0&0\\K&K&0\\K&K&K\end{pmatrix}.
\]
This is ordinary lower-triangular matrix multiplication, since every
nonzero coefficient transition and product is the indicated field map.
A $Q$-linear functional $Q\to Q$ is $q\mapsto qa$, and it is
admissible at $l$ exactly when $al\leq1$. Hence $\delta_\bot$ has
values $(\bot,u,u)$, $\delta_1$ is the identity, and $\delta_u$ is
zero. The support test holds already through the $l=1$ term, and
\[
 I_L=\begin{pmatrix}0&0&0\\K&K&0\\K&K&0\end{pmatrix}.
\]
The three diagonal simple types have their projective covers in the
three column ideals. Only the middle simple receives nonzero $I_L$
action. Its cover is the middle column, with groups $0,K,K$ and
identity transition $1\to u$. Thus the projectives with support $Q$
have precisely the groups $0,K^n,K^n$, for one integer $n\geq0$.
Every $Q$-linear support automorphism is multiplication by $a$, and
only $a=1$ is bijective; there is no further orbit identification.
The value $n=0$ retains the nontrivial support with zero amplitude,
and $n=1$ recovers the free rank-one semimodule.
\end{example}
